\documentclass[../../main.tex]{subfiles}
\begin{document}
{\bf [解]}

\begin{figure}[htb]
  \centering
  \begin{tikzpicture}[scale=2.0]
    % Triangle ABC
    \coordinate (A) at (60:3);
    \coordinate (B) at (0,0);
    \coordinate (C) at (3,0);

    % Points P, Q, R on sides
    \coordinate (P) at (0.8,0);
    \coordinate (Q) at ($(C)!0.267!(A)$);
    \coordinate (R) at ($(A)!0.267!(B)$);

    % Intersections
    \coordinate (A') at (intersection of A--P and C--R);
    \coordinate (B') at (intersection of A--P and B--Q);
    \coordinate (C') at (intersection of B--Q and C--R);

    % Draw main triangle
    \draw[thick] (A) -- (B) -- (C) -- cycle;

    % Draw lines AP, BQ, CR
    \draw (A) -- (P);
    \draw (B) -- (Q);
    \draw (C) -- (R);

    % Node labels
    \node[above] at (A) {$A$};
    \node[below left] at (B) {$B$};
    \node[below right] at (C) {$C$};

    \node[below] at (P) {$P$};
    \node[right] at (Q) {$Q$};
    \node[left] at (R) {$R$};

    \node[above right] at (A') {$A'$};
    \node[left] at (B') {$B'$};
    \node[below] at (C') {$C'$};

    % Length labels
    \node[below] at ($(B)!0.5!(P)$) {$x$};
    \node[right] at ($(C)!0.5!(Q)$) {$x$};
    \node[left] at ($(A)!0.5!(R)$) {$x$};

    \node[right] at ($(A)!0.5!(A')$) {$p$};
    \node[right] at ($(A')!0.5!(B')$) {$r$};
    \node[right] at ($(B')!0.5!(P)$) {$q$};

    \node[above] at ($(B)!0.5!(B')$) {$p$};
    \node[above] at ($(B')!0.5!(C')$) {$r$};
    \node[above] at ($(C')!0.5!(Q)$) {$q$};

    \node[left] at ($(C)!0.5!(C')$) {$p$};
    \node[left] at ($(C')!0.5!(A')$) {$r$};
    \node[left] at ($(A')!0.5!(R)$) {$q$};
  \end{tikzpicture}
  \caption{正三角形 $\triangle A'B'C'$ を作る補助線 $AP,BQ,CR$}
\end{figure}

題意の条件
\begin{align}
  0 < x < \dfrac{1}{2} \label{eq:1}
\end{align}
の元で考える．

\paragraph{(1)}

対称性から実数$p,q,r$を用いて
\begin{align}
  \begin{cases}
    |BB'| = |CC'| = |AA'| = p \\
    |PB'| = |QC'| = |RA'| = q \\
    |AB'| = |BC'| = |CA'| = r \\
    |BP | = |CQ|  = |AR|  = x
  \end{cases}\label{eq:2}
\end{align}
とおける．

\begin{figure}[htb]
  \centering
  % =========================================================
  % (a) 三角形 BC'R と直線 AP (点 B', A', A)
  % =========================================================
  \begin{tikzpicture}[scale=1.7, >=stealth, every node/.style={font=\small}]

    % 座標定義
    \coordinate (A) at (60:3);
    \coordinate (B) at (0,0);
    \coordinate (C) at (3,0);

    \coordinate (P) at (0.8,0);
    \coordinate (Q) at ($(C)!0.267!(A)$);
    \coordinate (R) at ($(A)!0.267!(B)$);

    \coordinate (A') at (intersection of A--P and C--R);
    \coordinate (B') at (intersection of A--P and B--Q);
    \coordinate (C') at (intersection of B--Q and C--R);

    % 対象三角形 \triangle BC'R の着色
    \fill[blue!12] (B) -- (C') -- (R) -- cycle;

    % \triangle BC'R の輪郭
    \draw[very thick, blue!85!black] (B) -- (C') -- (R) -- cycle;

    % 辺の延長部 (BR の延長上に A がある)
    \draw[dashed, gray!70] (R) -- (A);

    % メネラウス線 (直線 AP: 直線 A-A'-B'-P)
    \draw[line width=1.4pt, red!85!black] (A) -- (P);

    % 主要点のプロット
    \fill (B) circle (1.0pt) node[below left] {$B$};
    \fill (C') circle (1.0pt) node[below] {$C'$};
    \fill (R) circle (1.0pt) node[left] {$R$};
    \fill[red!85!black] (A) circle (1.0pt) node[above] {$A$};
    \fill[red!85!black] (B') circle (1.0pt) node[left] {$B'$};
    \fill[red!85!black] (A') circle (1.0pt) node[above right] {$A'$};

    % 長さ・比のラベル
    \node[blue!90!black, font=\scriptsize, above left] at ($(B)!0.5!(B')$) {$p$};
    \node[blue!90!black, font=\scriptsize, above left] at ($(B')!0.5!(C')$) {$r$};
    \node[teal!90!black, font=\scriptsize, left] at ($(C')!0.5!(A')$) {$r$};
    \node[teal!90!black, font=\scriptsize, left] at ($(A')!0.5!(R)$) {$q$};
    \node[purple!90!black, font=\scriptsize, left] at ($(A)!0.5!(R)$) {$x$};
    \node[purple!90!black, font=\scriptsize, left] at ($(B)!0.4!(A)$) {$1$};

    % タイトル・定理式
    \node[font=\footnotesize\bfseries, align=center] at (1.2, -0.6)
    {$\triangle BC'R$ と直線 $AP$\\[2pt]
    $\dfrac{B'C'}{BB'} \cdot \dfrac{A'R}{C'A'} \cdot \dfrac{AB}{RA} = 1$};

  \end{tikzpicture}
  \hspace{0.5cm}
  % =========================================================
  % (b) 三角形 APC と直線 BQ (点 B', Q, B)
  % =========================================================
  \begin{tikzpicture}[scale=1.7, >=stealth, every node/.style={font=\small}]

    % 座標定義
    \coordinate (A) at (60:3);
    \coordinate (B) at (0,0);
    \coordinate (C) at (3,0);

    \coordinate (P) at (0.8,0);
    \coordinate (Q) at ($(C)!0.267!(A)$);
    \coordinate (R) at ($(A)!0.267!(B)$);

    \coordinate (A') at (intersection of A--P and C--R);
    \coordinate (B') at (intersection of A--P and B--Q);
    \coordinate (C') at (intersection of B--Q and C--R);

    % 対象三角形 \triangle APC の着色
    \fill[orange!15] (A) -- (P) -- (C) -- cycle;

    % \triangle APC の輪郭
    \draw[very thick, orange!90!black] (A) -- (P) -- (C) -- cycle;

    % 辺の延長部 (CP の延長上に B がある)
    \draw[dashed, gray!70] (P) -- (B);

    % メネラウス線 (直線 BQ: 直線 B-B'-C'-Q)
    \draw[line width=1.4pt, red!85!black] (B) -- (Q);

    % 主要点のプロット
    \fill (A) circle (1.0pt) node[above] {$A$};
    \fill (P) circle (1.0pt) node[below] {$P$};
    \fill (C) circle (1.0pt) node[below right] {$C$};
    \fill[red!85!black] (B) circle (1.0pt) node[below left] {$B$};
    \fill[red!85!black] (B') circle (1.0pt) node[left] {$B'$};
    \fill[red!85!black] (Q) circle (1.0pt) node[right] {$Q$};

    % 長さ・比のラベル
    \node[orange!90!black, font=\scriptsize, right] at ($(A)!0.5!(B')$) {$p+r$};
    \node[orange!90!black, font=\scriptsize, right] at ($(B')!0.5!(P)$) {$q$};
    \node[teal!90!black, font=\scriptsize, right] at ($(C)!0.5!(Q)$) {$x$};
    \node[teal!90!black, font=\scriptsize, right] at ($(Q)!0.5!(A)$) {$1-x$};
    \node[purple!90!black, font=\scriptsize, below] at ($(B)!0.5!(P)$) {$x$};
    \node[purple!90!black, font=\scriptsize, below] at ($(B)!0.5!(C)$) {$1$};

    % タイトル・定理式
    \node[font=\footnotesize\bfseries, align=center] at (1.6, -0.6)
    {$\triangle APC$ と直線 $BQ$\\[2pt]
    $\dfrac{QA}{CQ} \cdot \dfrac{B'P}{AB'} \cdot \dfrac{BC}{PB} = 1$};

  \end{tikzpicture}
  \caption{メネラウスの定理を適用する2つの三角形と直線}
  \label{fig:menelaus_decomposition}
\end{figure}

三角形$BC'R, ABP$，および三角形$APC, BQC$にメネラウスの定理を用いて
\begin{align}
  & \dfrac{B'C'}{BB'} \cdot \dfrac{A'R}{C'A'} \cdot \dfrac{AB}{RA} = 1 \label{eq:2} \\
  & \dfrac{QA}{CQ} \cdot \dfrac{B'P}{AB'} \cdot \dfrac{BC}{PB} = 1 \label{eq:3}
\end{align}
とかける．ここに\cref{eq:2}を代入して
\begin{align}
  \begin{cases}
    \dfrac{r}{p} \cdot \dfrac{q}{r} \cdot \dfrac{1}{x} = 1 \\
    \dfrac{1-x}{x} \cdot \dfrac{q}{p+r} \cdot \dfrac{1}{x} = 1
  \end{cases}
  \therefore
  \begin{cases}
    r = \dfrac{1-2x}{x} p \\
    q = xp
  \end{cases} \label{eq:4}
\end{align}
である．

一方，$\triangle BCQ$ に余弦定理を用いて，
\begin{align}
  |BQ|^2    &= |BC|^2+|CQ|^2 - 2|BC|CQ|\cos\angle BCQ \\
  (p+q+r)^2 &= 1 + x^2 - 2x\cos\dfrac{\pi}{3} \\
  \therefore
  p + q + r &= \sqrt{x^2 - x + 1} \label{eq:5}
\end{align}
である．

\cref{eq:4}を\cref{eq:5}に代入して，
\begin{align}
  \Bigl( 1 + x + \dfrac{1-2x}{x} \Bigr) p = \sqrt{x^2 - x + 1}
  \therefore p = \dfrac{x}{\sqrt{x^2 - x + 1}} \label{eq:6}
\end{align}
となる．従って\cref{eq:4}に代入して
\begin{align}
  r = \dfrac{1-2x}{\sqrt{x^2 - x + 1}}, \quad q = \dfrac{x^2}{\sqrt{x^2 - x + 1}} \label{eq:7}
\end{align}
となる．
\cref{eq:6,eq:7}をまとめて，求める長さは
\begin{align}
  \begin{cases}
    |BB'| = p = \dfrac{x}{\sqrt{x^2 - x + 1}} \\
    |PB'| = q = \dfrac{x^2}{\sqrt{x^2 - x + 1}}
  \end{cases} \label{eq:8}
\end{align}
である．

\paragraph{(2)}

$\triangle A'B'C'$ は一辺の長さ $r$ の正三角形でその面積 $T$ は\cref{eq:7}より
\begin{align}
  T = \dfrac{\sqrt{3}}{4} r^2 = \dfrac{\sqrt{3}}{4} \cdot \dfrac{4x^2 - 4x + 1}{x^2 - x + 1}
\end{align}
と表せる．これが $\dfrac{1}{2} \cdot \dfrac{\sqrt{3}}{4}$ に等しいので
\begin{align}
  & \dfrac{1}{2} = \dfrac{4x^2 - 4x + 1}{x^2 - x + 1} \\
  & 8x^2 - 8x + 2 = x^2 - x + 1 \\
  & 7x^2 - 7x + 1 = 0 \\
  \therefore
  & x = \dfrac{7 \pm \sqrt{21}}{14}
\end{align}
である．
このうち\cref{eq:1}を満たすのは複合負で，求める$x$は
\begin{align}
  x = \dfrac{7 - \sqrt{21}}{14}
\end{align}
である．

\end{document}
